% ============================================================================
%  Lista de símbolos y abreviaturas
% ============================================================================
\chapter*{Lista de Símbolos y Abreviaturas}
\addcontentsline{toc}{chapter}{Lista de Símbolos y Abreviaturas}
\markboth{Lista de Símbolos y Abreviaturas}{Lista de Símbolos y Abreviaturas}

\section*{Siglas y abreviaturas}

\begin{longtable}{p{2.6cm}p{12.5cm}}
ADRC & Control con rechazo activo de perturbaciones (\textit{active disturbance rejection control}). \\
ARX & Modelo autorregresivo con entrada exógena (\textit{autoregressive with exogenous input}). \\
ASM1 & Modelo de lodos activados n.\textsuperscript{o}~1 de la IWA. \\
AW & Anti-\emph{windup}, manejo de la saturación del actuador. \\
BSM1 & \textit{Benchmark Simulation Model} n.\textsuperscript{o}~1 de la IWA. \\
DBO / DQO & Demanda bioquímica / química de oxígeno. \\
eADRC & ADRC formulado sobre el error de seguimiento (\textit{error-based ADRC}). \\
ESO & Observador de estado extendido (\textit{extended state observer}). \\
FOPTD & Primer orden con retardo (\textit{first order plus time delay}). \\
GPI & Control proporcional-integral generalizado. \\
HIL & \textit{Hardware in the loop}. \\
IAE / ISE / ITAE & Integral del error absoluto / cuadrático / absoluto ponderado por el tiempo. \\
IWA & \textit{International Water Association}. \\
MAD & Desviación absoluta mediana (\textit{median absolute deviation}). \\
MPC & Control predictivo basado en modelo. \\
OD & Oxígeno disuelto. \\
OE1--OE5 & Objetivos específicos 1 a 5 de este trabajo. \\
OU & Proceso de Ornstein-Uhlenbeck. \\
PID / PI & Controlador proporcional-integral-derivativo / proporcional-integral. \\
PLC & Controlador lógico programable. \\
PTAR & Planta de tratamiento de aguas residuales. \\
SCADA & Sistema de supervisión, control y adquisición de datos. \\
ST & Texto estructurado (lenguaje de la norma IEC 61131-3). \\
UASB & Reactor anaerobio de flujo ascendente (\textit{upflow anaerobic sludge blanket}). \\
ZOH & Retenedor de orden cero (\textit{zero-order hold}). \\
\end{longtable}

\section*{Símbolos del proceso}

\begin{longtable}{p{2.6cm}p{12.5cm}}
$S_{O,i}$ & Concentración de oxígeno disuelto en el reactor $i$ [mg/L o g/m$^3$]. \\
$S_O^{*}$ & Concentración de saturación de oxígeno; 8~g/m$^3$ en el BSM1. \\
$\KLa$ & Coeficiente de transferencia de oxígeno del reactor 5 [d$^{-1}$]; variable manipulada en el BSM1. \\
$r_{Z,5}$ & Tasa neta de respiración biológica en el reactor 5 [g/(m$^3\cdot$d)]. \\
$Q_0,\,Q_a,\,Q_r,\,Q_w$ & Caudales de entrada, recirculación interna, recirculación de lodos y purga. \\
$V_i$ & Volumen del reactor $i$ [m$^3$]. \\
$X_{B,H},\,X_{B,A}$ & Biomasa heterótrofa y autótrofa activa [g~DQO/m$^3$]. \\
$S_S,\,S_{NH}$ & Sustrato rápidamente biodegradable y nitrógeno amoniacal. \\
$k$ & Constante del sumidero de oxígeno de primer orden, $r = k\,S_O$ [h$^{-1}$]. \\
\end{longtable}

\section*{Símbolos del control}

\begin{longtable}{p{2.6cm}p{12.5cm}}
$y$ & Salida medida: oxígeno disuelto según el sensor instalado [mg/L]. \\
$r$ & Referencia de oxígeno disuelto [mg/L]. \\
$e$ & Error de seguimiento; en el texto $e_y = y-r$ y $e = r-y$ se distinguen explícitamente. \\
$u$ & Señal de control en unidades del PLC [\% de consigna del soplador]. \\
$u_0$ & Señal de control virtual, antes de dividir por $\bo$. \\
$\varepsilon$ & Perturbación a la entrada de la planta, en las mismas unidades de $u$ [\%]. \\
$\eta$ & Ruido de medida del sensor de OD, sumado a la salida [mg/L]. \\
$\bar{\varepsilon}$ & Perturbación equivalente a la entrada: $\varepsilon$ más la dinámica de la planta que no es un integrador; en primer orden, $\bar{\varepsilon} = \varepsilon - y/k_{dc}$ [\%]. \\
$\bar{\xi}$ & Perturbación total sobre la salida: todo lo que mueve $\dot y$ y no es $\bo u$, $\bar{\xi} = \bo\,\bar{\varepsilon}$ [mg/(L$\cdot$h)]. \\
$\xi$ & Perturbación total vista desde el error, $\xi = \bar{\xi} - u^*$; con referencia constante, $\xi = \bar{\xi}$ [mg/(L$\cdot$h)]. \\
$\hat{\xi},\ \hat{y}$ & Estimaciones del observador de la perturbación total y de la salida. \\
$\bo$ & Ganancia de alta frecuencia del modelo $\dot y=\bar{\xi}+\bo u$ [mg/L/(\%$\cdot$h)]. \\
$K$ & Ganancia estática incremental [mg/L por punto de señal de control]. \\
$\tau$ & Constante de tiempo del lazo abierto [min]. \\
$\theta$ & Retardo de transporte y de medida [s]. \\
$\wc$ & Ancho de banda del lazo cerrado (polo del controlador) [rad/s o h$^{-1}$]. \\
$\wo$ & Ancho de banda del observador [rad/s, h$^{-1}$ o d$^{-1}$]. \\
$L=[l_1\ l_2]^{\!\top}$ & Ganancias del observador de estado extendido. \\
$k_0,k_1,k_2$ & Ganancias del controlador GPI. \\
$K_P,\,T_N,\,T_V,\,T_f$ & Ganancia, tiempo integral, tiempo derivativo y filtro de la derivada del PID. \\
$E_1(s)$ & Polinomio mónico estable del esquema de anti-\emph{windup}. \\
$\Ts$ & Periodo de muestreo [s]. \\
$\sigma(u)$ & Desviación típica de la señal de control; medida de esfuerzo. \\
GM, PM & Margen de ganancia [dB] y margen de fase [grados]. \\
\end{longtable}
